\documentclass[10pt,letterpaper]{article}
\usepackage[margin=1in]{geometry}
\usepackage[utf8]{inputenc}
\usepackage[T1]{fontenc}
\usepackage{txfonts}
\usepackage{microtype}
\usepackage{amsmath}
\usepackage{amssymb}
\usepackage{booktabs}
\usepackage{array}
\usepackage{enumitem}
\usepackage{xcolor}
\usepackage{hyperref}
\usepackage{titlesec}
\usepackage{abstract}
\usepackage{authblk}
\usepackage{caption}
\usepackage{graphicx}

\hypersetup{
  colorlinks=true,
  linkcolor=black,
  urlcolor=blue!50!black,
  citecolor=blue!50!black,
}

\titleformat{\section}{\normalfont\large\bfseries}{\thesection}{0.6em}{}
\titleformat{\subsection}{\normalfont\normalsize\bfseries}{\thesubsection}{0.6em}{}

\setlength{\parskip}{0.4em}
\setlength{\parindent}{0pt}
\frenchspacing

\title{\bfseries Cheap PRMs:\\
Multi-Dimensional Process Reward Modeling for\\
Domain-Specialized Reasoning}

\author{Rama Krishna Bachu\\
\small Bottensor, Independent Research\\
\small \texttt{ramakrishna.bachu@bottensor.xyz}}

\date{April 2026}

\begin{document}

\maketitle

\begin{abstract}
\noindent
Process Reward Models (PRMs) score individual reasoning steps for correctness and have become central to recent reasoning-model training pipelines. Most published PRMs are math-only and use a single scalar correctness signal; we examine whether multi-dimensional scoring on a domain-specialized reasoning task (DeFi/crypto financial analysis) is worth the extra modeling surface.

We describe NPC Fin-PRM 7B, a 7B QLoRA-trained PRM scored on four dimensions --- \texttt{factual\_accuracy}, \texttt{logical\_validity}, \texttt{completeness}, and \texttt{risk\_awareness} --- trained in 17.4 hours on a single H100 from approximately 80{,}000 step-level judge labels generated by Qwen2.5-72B over 4{,}866 reasoning trees.

On a stratified 200-example held-out validation split, the model achieves Spearman 0.92 against judge labels (rating accuracy 88.5\%, error-detection F1 0.84) at MAE 0.04 on the 0--1 score scale. Out-of-distribution evaluation on 307 gold-correct math-reasoning steps from GSM8K and MATH-500 finds only \textbf{5.2\%} mis-flagged as flawed and a mean \texttt{overall\_score} of 0.856 --- substantially better cross-domain transfer than expected for a DeFi-only training corpus, with a curious side-effect that the model extrapolates beyond its training labels by emitting \texttt{EXCELLENT} and \texttt{PERFECT} ratings on 3.9\% of OOD math steps despite never being trained to produce them.

Two findings stand out. First, three of the four dimensions form a tightly correlated cluster (pairwise Spearman 0.85--0.92) that is largely captured by a single axis; the judge's overall score is 95\% explained by \texttt{logical\_validity} alone. Second, despite Spearman 0.92 the model is poorly calibrated as a probability (ECE = 0.21): it over-flags in the 0.1--0.5 score band and under-flags around 0.5--0.7, so the score scale should not be used as a calibrated probability without a Platt-scaling or isotonic-regression sidecar.

The contribution is recipe-level: a complete pipeline for cheap, domain-specialized process reward modeling on accessible hardware, with the limitations and unmet experiments reported alongside the wins.
\end{abstract}

\section{Introduction}

A class of recent work argues that scaling test-time reasoning depth can substitute for scaling pre-training compute \cite{lightman2023verify, deepseek2025r1}. Process Reward Models (PRMs) are the supervision instrument behind this line of work: instead of scoring an entire chain-of-thought against a final answer, a PRM scores each \emph{step} of reasoning, providing a much denser feedback signal that can be used for best-of-$N$ reranking, search-guided decoding, or as the reward in reinforcement-learning fine-tuning.

The published PRM literature is small but growing \cite{wang2024mathshepherd, wang2024openr, luo2024omegaprm}, and it is dominated by two patterns. First, almost every published PRM is trained for mathematical reasoning, on derivatives of GSM8K or MATH; few address domain-specialized reasoning where the failure modes look qualitatively different (legal, medical, financial). Second, almost every published PRM is trained as a binary classifier (\texttt{step correct} vs.\ \texttt{step incorrect}); few decompose the correctness signal into orthogonal axes that distinguish, say, factual error from logical fallacy from incomplete consideration of risk.

This paper documents one attempt at both choices: a 7B PRM for DeFi/crypto financial reasoning, scored on four dimensions, trained in 17.4 hours of single-H100 wall-clock time with all training labels generated by an open-weight judge model (Qwen2.5-72B) served on the same hardware. We are interested in three questions:
\begin{enumerate}[leftmargin=*]
    \item Does training on multi-dimensional judge labels yield a usable PRM for a non-math domain?
    \item Are the four dimensions actually independent, or do they collapse?
    \item How well does the resulting model rank steps, and how well does its score scale function as a probability of error?
\end{enumerate}

We answer all three quantitatively. The model ranks steps well (Spearman 0.92 on a held-out test split), the four dimensions are not independent (three cluster at $\rho \in [0.85, 0.92]$), and the score scale is miscalibrated as a probability (ECE 0.21) despite the strong rank correlation. We report the recipe and the limitations in the same document as the wins. The contribution is recipe-level, not method-level.

The model is publicly available: \href{https://huggingface.co/ramankrishna10/npc-fin-prm-7b}{\texttt{ramankrishna10/npc-fin-prm-7b}}. Training scripts, eval harness, and analysis scripts are at \href{https://github.com/ramankrishna/bottensor-models/tree/main/training/npc-fin-prm-7b}{\texttt{bottensor-models}}.

\section{Related Work and Positioning}
\label{sec:related}

The closest prior PRMs are Math-Shepherd \cite{wang2024mathshepherd}, OpenR \cite{wang2024openr}, OmegaPRM \cite{luo2024omegaprm}, and the original PRM800K work \cite{lightman2023verify}. All four are trained for math reasoning; all four predict a single scalar reward per step. The labeling protocols differ --- PRM800K used human raters; Math-Shepherd and OmegaPRM use automatic process supervision via Monte Carlo rollouts; OpenR is a framework that supports both --- but the design space they explore is the math-correctness scalar.

DeepSeek-R1 \cite{deepseek2025r1} reports using process-style supervision but does not release a separable PRM artifact. The OpenAI o1 line acknowledges process supervision but is closed.

Our positioning is along two axes. First, we move the domain from math to DeFi/crypto financial reasoning, where the failure modes are dominated by \texttt{LOGICAL\_FALLACY} and \texttt{OMISSION} (collectively 92.8\% of negative labels in our data) rather than the arithmetic and algebraic errors of GSM8K. Second, we decompose the reward signal into four named dimensions and report on whether the decomposition is empirically useful or whether it collapses to one effective axis.

The cheap-recipe positioning of this paper is shared with the NPC Fast preprint \cite{bachu2026npcfast}, which trains a 1.7B router on a single H100 with a similar accessible-hardware focus. Our contribution here is the same kind of recipe-level documentation in the PRM space.

\section{Data Pipeline}
\label{sec:data}

\subsection{Logic trees}

We start from a corpus of 4{,}866 reasoning ``logic trees,'' each consisting of a DeFi/crypto scenario followed by 1--7 reasoning steps that collectively analyze the scenario (token price moves, on-chain events, exploit hypotheses, treasury decisions, etc.). The tree corpus is internal and not redistributed.

Each step within each tree becomes one labeled training example. With an average tree depth of approximately 16 reasoning steps when including alternative branches, the corpus yields approximately 80{,}000 labelable step records. We retain a step record only if (a) the judge output parses as valid JSON against the schema, (b) the full prompt+context fits within the 2{,}048-token training context, (c) the reasoning chain it belongs to passes a MinHash near-duplicate filter against the rest of the corpus, and (d) the chain meets minimum quality heuristics ($\geq$3 reasoning steps, contains specific metrics or levels, internally consistent, no hallucinated token names). The exact per-filter drop counts were not preserved in the training metadata; the cumulative effect reduces 80{,}000 candidate step records to 41{,}713 retained examples (37{,}542 train, 4{,}171 validation), a 47.9\% retention rate.

\subsection{Judge labeling}

For each step, we present its surrounding context (scenario plus prior steps) to a judge model and ask for a structured JSON evaluation:

\begin{quote}\small\ttfamily
\{ "overall\_score": 0.85,\\
\hspace*{1em}"rating": "STRONG | ACCEPTABLE\_WITH\_ISSUES | FLAWED",\\
\hspace*{1em}"dimensions": \{\\
\hspace*{2em}"factual\_accuracy":  \{"score": 0.95, "justification": "..."\},\\
\hspace*{2em}"logical\_validity":  \{"score": 0.85, "justification": "..."\},\\
\hspace*{2em}"completeness":      \{"score": 0.80, "justification": "..."\},\\
\hspace*{2em}"risk\_awareness":   \{"score": 0.75, "justification": "..."\}\\
\hspace*{1em}\},\\
\hspace*{1em}"explanation": "...",\\
\hspace*{1em}"error\_identified": "LOGICAL\_FALLACY | OMISSION | UNJUSTIFIED\_CONFIDENCE | FACTUAL\_ERROR | LOGICAL\_LEAP | null"\\
\}
\end{quote}

The judge is Qwen2.5-72B-Instruct, served locally on the same H100 via vLLM, amortizing the rented compute across both jobs. The labeling prompt instructs the judge to score each dimension independently on a 0--1 scale, with a brief justification per dimension and a single ``error\_identified'' tag drawn from a closed set when the rating is below \texttt{STRONG}.

We confirm post-hoc that the JSON parse-failure rate on the held-out validation set is exactly 0 of 4{,}171 records (Section~\ref{sec:eval}). The judge is robust at this output schema.

\subsection{Label distribution}

The 4{,}171 validation records distribute across rating categories as 55.6\% \texttt{STRONG}, 30.1\% \texttt{ACCEPTABLE\_WITH\_ISSUES}, and 14.3\% \texttt{FLAWED}. The mean overall\_score is 0.726; the distribution is bimodal with peaks at 0.6--0.7 ($n=555$) and 0.8--0.9 ($n=2{,}121$). Step-depth (number of prior steps in context) is approximately uniform across depths 0--3 (920--950 each) and tapers at depth 4+ (412 at depth 4, 67 at depth 5).

The error-tag distribution is dominated by two failure modes: \texttt{LOGICAL\_FALLACY} (57.2\% of non-null tags) and \texttt{OMISSION} (35.6\%). \texttt{UNJUSTIFIED\_CONFIDENCE} (4.5\%), \texttt{FACTUAL\_ERROR} (2.6\%), and \texttt{LOGICAL\_LEAP} (0.1\%) make up the long tail. We note for the paper that the 92.8\% mass on two tags is a finding of its own: future PRMs in this domain may want to expand the failure taxonomy.

\subsection{Are the four dimensions independent?}
\label{sec:orthogonality}

Pairwise Spearman correlations among the four dimensions in the gold labels (Table~\ref{tab:orth}) reveal that three of them form a tightly correlated cluster.

\begin{table}[!ht]
\centering
\small
\caption{Pairwise Spearman correlation among the four dimensions in the gold labels ($n = 4{,}171$).}
\label{tab:orth}
\begin{tabular}{lrrrr}
\toprule
 & factual & logical & complete & risk \\
\midrule
factual\_accuracy   & 1.000 & 0.743 & 0.772 & 0.715 \\
logical\_validity   &       & 1.000 & 0.851 & 0.916 \\
completeness        &       &       & 1.000 & 0.863 \\
risk\_awareness     &       &       &       & 1.000 \\
\bottomrule
\end{tabular}
\end{table}

\texttt{logical\_validity} $\leftrightarrow$ \texttt{risk\_awareness} correlate at $\rho = 0.916$. The \texttt{logical} / \texttt{completeness} / \texttt{risk} triangle is at $\rho \in [0.85, 0.92]$; \texttt{factual\_accuracy} is meaningfully more independent at $\rho \in [0.71, 0.77]$.

A complementary view: regress the judge's \texttt{overall\_score} against each individual dimension (Table~\ref{tab:explain}). The judge's overall score is 95\% explained by \texttt{logical\_validity} alone.

\begin{table}[!ht]
\centering
\small
\caption{Spearman correlation between the judge's \texttt{overall\_score} and each dimension.}
\label{tab:explain}
\begin{tabular}{lr}
\toprule
Dimension & $\rho$ vs.\ overall \\
\midrule
\texttt{logical\_validity}   & 0.955 \\
\texttt{risk\_awareness}     & 0.949 \\
\texttt{completeness}        & 0.864 \\
\texttt{factual\_accuracy}   & 0.762 \\
\bottomrule
\end{tabular}
\end{table}

This is the paper's first headline finding: \emph{the four-dimensional decomposition is partially redundant}. Future work in PRM design should consider more orthogonal axes (e.g.\ factual\_accuracy + logical\_validity + uncertainty\_calibration + counterfactual\_robustness) rather than the four we used. The result also bounds the practical contribution of any per-dimension downstream pipeline: a single \texttt{logical\_validity} head would recover almost all of the \texttt{overall\_score} signal.

\section{Architecture and Training}
\label{sec:training}

We fine-tune \texttt{Qwen/Qwen2.5-7B-Instruct} via QLoRA on the labeled step data, treating the task as a chat-style conditional-generation problem (predict the JSON evaluation given the scenario and step). Training configuration is in Table~\ref{tab:train}.

\begin{table}[!ht]
\centering
\small
\caption{Training configuration. Numbers verified from the published \texttt{trainer\_state.json} + \texttt{training\_metadata.json}.}
\label{tab:train}
\begin{tabular}{ll}
\toprule
\textbf{Field} & \textbf{Value} \\
\midrule
Base model           & \texttt{Qwen/Qwen2.5-7B-Instruct} \\
Method               & QLoRA, 4-bit NF4 base + bf16 LoRA adapters \\
LoRA                 & $r=32$, $\alpha=64$, dropout $0.05$ \\
Target modules       & \texttt{q,k,v,o,gate,up,down\_proj} (all 7) \\
Optimizer            & 8-bit AdamW, paged \\
LR schedule          & cosine, peak $2 \times 10^{-4}$, warmup 0.05 \\
Effective batch      & 16 (per-device 2 $\times$ grad-accum 8) \\
Max sequence length  & 2{,}048 tokens \\
Total steps          & 7{,}041 (3 epochs $\times$ 2{,}346 steps/epoch) \\
Train examples       & 37{,}542 \\
Val examples         & 4{,}171 \\
Final train loss     & 0.2518 \\
Final eval loss      & 0.2168 \\
Wall clock           & 17.4 hours \\
Hardware             & Single NVIDIA H100 80\,GB \\
Peak VRAM            & 23.69\,GB \\
\bottomrule
\end{tabular}
\end{table}

The training run completed cleanly: eval loss descended monotonically from 0.4839 at step 200 to 0.2158 at step 6{,}200 with no LR resets, no loss spikes, and no detectable overfitting on the validation split. The full 35-eval-point loss curve is in the supplementary material.

\paragraph{Note on the bootstrap claim.} The published model card describes the training as ``3 bootstrapped iterations: train, evaluate, re-label worst 500, retrain.'' Inspection of the trainer state, however, shows a single continuous training run of 7{,}041 steps at $\lfloor 7041 / 3 \rfloor = 2{,}347 \approx$ steps-per-epoch, with no detectable LR resets or loss spikes that a per-iteration restart would produce. The model card description was inherited from an earlier project plan; the actual run was three epochs of vanilla SFT on a fixed label set. We document this honestly here and treat the iterated relabeling as future work in Section~\ref{sec:limitations}.

\section{Evaluation}
\label{sec:eval}

\subsection{Held-out test set}

We evaluate on a stratified 200-example sample of the validation split (80 \texttt{STRONG} / 80 \texttt{ACCEPTABLE} / 40 \texttt{FLAWED}). The PRM is run in greedy decoding with a 768-token output budget; we use a tolerant JSON parser that recovers scores from truncated JSON via regex when the budget is exhausted mid-justification. The parse-success rate is 100\%.

\begin{table}[!ht]
\centering
\small
\caption{In-distribution evaluation on a stratified 200-example test split. ``Card claims'' are reproduced from the published model card and confirmed by this independent re-run.}
\label{tab:eval-id}
\begin{tabular}{lrr}
\toprule
\textbf{Metric} & \textbf{Card claim} & \textbf{This work ($n$=200)} \\
\midrule
Spearman, overall\_score                & 0.94    & \textbf{0.9234} \\
Step-level rating accuracy              & 89.2\%  & \textbf{88.5\%} \\
Error-detection F1 (\texttt{FLAWED})    & 0.87    & \textbf{0.8421} \\
MAE on overall\_score                   & ---     & \textbf{0.0404} \\
JSON parse failures                     & ---     & \textbf{0 / 200} \\
\midrule
Spearman, factual\_accuracy             & ---     & 0.842 \\
Spearman, logical\_validity             & ---     & \textbf{0.931} \\
Spearman, completeness                  & ---     & 0.865 \\
Spearman, risk\_awareness               & ---     & 0.908 \\
\bottomrule
\end{tabular}
\end{table}

\begin{table}[!ht]
\centering
\small
\caption{Confusion matrix on the 200-example test split. Rows are gold ratings, columns are PRM predictions.}
\label{tab:confusion}
\begin{tabular}{lrrr}
\toprule
 & STRONG & ACCEPTABLE & FLAWED \\
\midrule
STRONG (n=80)         & \textbf{79} & 1  & 0  \\
ACCEPTABLE (n=80)     & 10          & \textbf{66} & 4  \\
FLAWED (n=40)         & 0           & 8  & \textbf{32} \\
\bottomrule
\end{tabular}
\end{table}

The PRM ranks steps well and rarely confuses extremes: it never predicts \texttt{FLAWED} for a gold-\texttt{STRONG} step, and never predicts \texttt{STRONG} for a gold-\texttt{FLAWED} step. The off-diagonal mass concentrates on adjacent-category confusions (\texttt{STRONG}$\leftrightarrow$\texttt{ACCEPTABLE} and \texttt{ACCEPTABLE}$\leftrightarrow$\texttt{FLAWED}), which is the expected failure mode of a graded classifier.

\subsection{Calibration}
\label{sec:calibration}

Despite the strong rank correlation, the score scale is miscalibrated as a probability of error. Treating $1 - \texttt{overall\_score}$ as predicted $P(\text{flawed})$ and binning into 10 deciles, we obtain ECE = 0.21 (Table~\ref{tab:calib}).

\begin{table}[!ht]
\centering
\small
\caption{Reliability table for $1 - \texttt{overall\_score}$ as predicted probability of \texttt{FLAWED}, on the 200-example test split. Empty bins omitted.}
\label{tab:calib}
\begin{tabular}{lrrr}
\toprule
$P(\text{flawed})$ bin & $n$ & avg.\ predicted & actual flawed rate \\
\midrule
$[0.0, 0.1)$ & 5 & 0.095 & 0.000 \\
$[0.1, 0.2)$ & 84 & 0.146 & 0.000 \\
$[0.2, 0.3)$ & 14 & 0.276 & 0.071 \\
$[0.3, 0.4)$ & 38 & 0.331 & 0.053 \\
$[0.4, 0.5)$ & 23 & 0.455 & 0.217 \\
$[0.5, 0.6)$ & 20 & 0.566 & \textbf{0.900} \\
$[0.6, 0.7)$ & 10 & 0.610 & 0.800 \\
$[0.7, 0.8)$ & 5 & 0.710 & 1.000 \\
$[0.8, 0.9)$ & 1 & 0.830 & 1.000 \\
\bottomrule
\end{tabular}
\end{table}

The PRM \emph{over-flags} in the 0.1--0.4 score band (predicted $\sim$0.15--0.33 flaw probability, actual $\sim$0.00--0.07) and \emph{under-flags} sharply in the 0.5--0.6 band (predicted 0.57, actual 0.90). At the extremes ($\geq 0.7$ predicted) calibration is good. The pattern is consistent with a model trained to imitate a labeler with a hard rating threshold (Section~\ref{sec:related}): it learns the ordering well but compresses probability mass in non-extreme regions.

\paragraph{Practical implication.} Downstream pipelines should not use $1 - \texttt{overall\_score}$ as a calibrated $P(\text{flawed})$. A standard fix is to fit a Platt-scaling or isotonic-regression layer on a held-out calibration split and apply it post-hoc; we leave the calibrated checkpoint to future work. For applications that only require ranking (best-of-$N$ reranking, search-guided decoding), the raw scores are usable as-is.

\subsection{Out-of-distribution behavior}

We probe the cost of domain specialization by running the PRM on 307 reasoning steps drawn from gold-correct GSM8K (134 steps) and MATH-500 (173 steps) solutions. Every step in this set is, by construction, a correct step in a correct mathematical derivation: a perfectly transferable PRM should rate every step \texttt{STRONG}. We sentence-tokenize the published gold solutions, present each tokenized step to the PRM with the surrounding earlier-step context, and parse the JSON output as in Section~\ref{sec:eval}.

The PRM mis-flags only \textbf{5.2\%} of these gold-correct math steps as \texttt{FLAWED} (16 of 307), with mean \texttt{overall\_score} of 0.856 --- substantially \emph{higher} than the in-distribution mean of 0.726. The full breakdown is in Table~\ref{tab:ood}; the parse-success rate is again 100\% (0/307 parse failures).

\begin{table}[!ht]
\centering
\small
\caption{Out-of-distribution evaluation on gold-correct math reasoning. ``mis-flagged'' = predicted rating \texttt{FLAWED}.}
\label{tab:ood}
\begin{tabular}{lrrrr}
\toprule
Source & $n$ & mean overall\_score & FLAWED & FLAWED \% \\
\midrule
GSM8K          & 134 & 0.869 & 6  & 4.5\% \\
MATH-500       & 173 & 0.846 & 10 & 5.8\% \\
\midrule
\textbf{Total} & 307 & 0.856 & 16 & 5.2\% \\
\bottomrule
\end{tabular}
\end{table}

This result is more positive for the PRM than we expected. A DeFi-trained step verifier rejects approximately one in twenty gold-correct math reasoning steps and gives a higher average score than it does on its in-distribution training domain. The 5.2\% rejection rate is small enough that, for a downstream pipeline that uses the PRM as a soft filter rather than a hard gate, naive cross-domain application is usable.

\paragraph{Rating extrapolation as a side effect.} The training rating set is \{\texttt{STRONG}, \texttt{ACCEPTABLE\_WITH\_ISSUES}, \texttt{FLAWED}\}. On the OOD math reasoning, however, the model emits \texttt{EXCELLENT} on 11 steps (3.6\%) and \texttt{PERFECT} on 1 step (0.3\%), labels it was never trained to produce. We attribute this to the underlying base model's prior over generic ``goodness'' lexicon plus the small temperature of 0.1: the model finds gold-correct math reasoning to be ``better than its training STRONG examples'' and reaches for vocabulary outside the trained schema. Downstream consumers should fold these into \texttt{STRONG} or impose a closed-set parser; we report the phenomenon because it is a real failure of schema discipline that an HF-published model card does not mention.

The full predicted-rating distribution on the 307-step OOD set is \texttt{STRONG} 87.9\% / \texttt{FLAWED} 5.2\% / \texttt{EXCELLENT} 3.6\% / \texttt{ACCEPTABLE\_WITH\_ISSUES} 2.9\% / \texttt{PERFECT} 0.3\%.

\section{Limitations}
\label{sec:limitations}

\subsection{Bootstrap loop claimed but not run}

As discussed in Section~\ref{sec:training}, the published model card describes a 3-iteration bootstrap with worst-500 relabeling that we could not validate from the trainer state. We document this here rather than hide it. Running the bootstrap properly --- training iter-1, evaluating, relabeling the worst Spearman residuals, retraining --- would add a real ablation to the paper. We leave this as future work.

\subsection{Judge ceiling}

All training labels come from a single judge model, Qwen2.5-72B-Instruct. The Spearman 0.92 number we report is correlation \emph{vs. the same family of judge}, which is bounded above by the judge's own quality. A genuine ceiling check would require human-labeled gold on a 200-step subset and Spearman against the human; we do not have that data.

\subsection{Dimensions partially redundant}

Section~\ref{sec:orthogonality} reports that three of the four dimensions cluster at $\rho \in [0.85, 0.92]$ in the gold labels. The four-dimension framing is therefore partially redundant. Future PRM designs in this space should pick more orthogonal axes.

\subsection{Single-domain}

The training corpus is DeFi/crypto financial reasoning. The OOD probe in the previous section quantifies the transfer cost to gold-correct math reasoning; we do not characterize transfer to other reasoning domains (legal, medical, scientific) and users should expect substantial failure on those.

\subsection{Calibration not corrected}

The shipped model is uncalibrated as a probability (Section~\ref{sec:calibration}). A Platt-scaling sidecar would correct this in approximately one minute of compute; we have not shipped that yet.

\subsection{No human-eval baseline}

We rely on judge-model labels throughout. A human-eval pass on a 100--200-step subset is the missing eval that would close the judge-ceiling gap.

\section{Conclusion}

We trained a 7B Process Reward Model for DeFi/crypto reasoning in 17.4 hours of single-H100 wall clock, using only judge-model labels from Qwen2.5-72B over 4{,}866 reasoning trees. On a stratified 200-example held-out test, the model achieves Spearman 0.92 against the judge labels and matches the published card's headline numbers. We report two findings the model card does not: the four labeled dimensions are not independent (three cluster at $\rho \in [0.85, 0.92]$), and the score scale is miscalibrated (ECE 0.21) despite the strong rank correlation.

The contribution is recipe-level. We document a complete pipeline for building a domain-specialized PRM on accessible hardware, with the unmet experiments reported alongside the wins. We share this analysis in the belief that recipes of this kind are undersupplied in the PRM literature, and that the cost of describing the failures is small relative to the value of letting other practitioners reproduce or improve on the result.

\subsection*{Code and Artifacts}

\begin{itemize}[leftmargin=*,itemsep=0pt]
\item Model: \href{https://huggingface.co/ramankrishna10/npc-fin-prm-7b}{\texttt{huggingface.co/ramankrishna10/npc-fin-prm-7b}}
\item Eval harness, analysis, paper-ready numbers: \href{https://github.com/ramankrishna/bottensor-models/tree/main/training/npc-fin-prm-7b}{\texttt{github.com/ramankrishna/bottensor-models}}
\end{itemize}

\subsection*{Acknowledgements}

This work was conducted as an independent open-source effort. Compute for both training and judge labeling was rented from Hyperbolic (single H100); evaluation was performed on a 24\,GB Apple M5 laptop using MLX 4-bit quantization. Training stack built on open-source contributions from the Qwen, Hugging Face TRL, PEFT, bitsandbytes, vLLM, and Unsloth projects.

\begin{thebibliography}{9}

\bibitem{lightman2023verify}
Lightman, H., Kosaraju, V., Burda, Y., Edwards, H., Baker, B., Lee, T., Leike, J., Schulman, J., Sutskever, I., \& Cobbe, K. (2023).
\textit{Let's Verify Step by Step.}
arXiv:2305.20050.

\bibitem{wang2024mathshepherd}
Wang, P., Li, L., Shao, Z., Xu, R., Dai, D., Li, Y., Chen, D., Wu, Y., \& Sui, Z. (2024).
\textit{Math-Shepherd: Verify and Reinforce LLMs Step-by-step without Human Annotations.}
arXiv:2312.08935.

\bibitem{wang2024openr}
Wang, J., Fang, M., Wan, Z., et al.\ (2024).
\textit{OpenR: An Open Source Framework for Advanced Reasoning with Large Language Models.}
arXiv:2410.09671.

\bibitem{luo2024omegaprm}
Luo, L., Liu, Y., Liu, R., et al.\ (2024).
\textit{Improve Mathematical Reasoning in Language Models by Automated Process Supervision.}
arXiv:2406.06592.

\bibitem{deepseek2025r1}
DeepSeek-AI. (2025).
\textit{DeepSeek-R1: Incentivizing Reasoning Capability in LLMs via Reinforcement Learning.}
arXiv:2501.12948.

\bibitem{uesato2022solving}
Uesato, J., Kushman, N., Kumar, R., Song, F., Siegel, N., \& Wang, L. (2022).
\textit{Solving math word problems with process- and outcome-based feedback.}
arXiv:2211.14275.

\bibitem{bachu2026npcfast}
Bachu, R. K. (2026).
\textit{NPC Fast 1.7B: Building a Usable Small Model on a Single H100.}
Zenodo. \href{https://doi.org/10.5281/zenodo.19771040}{doi.org/10.5281/zenodo.19771040}.

\end{thebibliography}

\end{document}
